% Page 070

\seite{70}

erhoben, zum Religionsunterricht unfähige Pfarrer immer heger und ob\ohb
rüchiger fahre und andrerseits auf offener Seite für die Großväterliche\unsicher\ob
aufsteht, daß die Sittenpolizei und die Kirchspielen so weit gehen, daß eine\ob
gute nützl. Vorstellung gestört worden. Man hoffte aber, daß die Spannung

31.\ März   Auf der Hälfte worden fünf entlassen 6 einheimische\ob
und 3 hies. Bickinder

Mit Rücksicht auf den fälligwerdenden Ostertermin beginnt\ob
der neue Kursus d.\ J. ausnahmsweise erst einmal mit\ob
d.\ 1. April. Die kommenden Tage finden also auf die Eün\ohb
weisung den Osterdienstags statt, wenn es d.\ J. aufgenom\ohb
men in vorklassenen 3 Mädchen sind.

15.\ April. Die sechswöchigen Osterferien beginnen mit heutigem\ob
Tage und enden mit dem 28.\ April.

Da der Tag des Ewigen Gebetes dieses Jahre in die\ob
Pfingstoktave (Pfingstdienstag den 7.\ Juni) fällt, wird\ob
dieser geschoben und die Osterferien gesetzt und\ob
für Beginn der Herbstferien wurde am 30.\ April.

6.\ Mai.    Ein bedeutungsvoller Tag liegt hinter uns – ein\ob
Großkampftag, der 4.\ Mai – der Tag der Erteilung\ob
der Bürgermeistereiwahls- u.\ Gemeinderatswahl.

Alle warten überall die Ergebnisse nächst- vom\ob
Rundfunk Strichberichten die Frankfurt und Berlin\ob
Wahlberechtigt in unserem Dorfe waren zur Reichstags\ohb
wahrwerk    zur Gemeinderatswahl u.\ Bürgermeistereirats\ohb
wahlerech.    Für die Reichstagswahl worden\ob
14 2 Stimmen abgegeben, davon entfielen\ob
auf das Zentrum 117, auf die christl. soziale Volksge\ohb
meinschaft 2, ungültig waren 3.\ Zur Gemeinderats\ohb
erteilwahl worden aufgestellt     Stimmen\ob
abgegeben. Auf den Wahlvorschlag Kasbach entfielen

\pend
\begin{verse}
Stimmen und den Wahlvorschlag Keil- Sefferw.\\
Stimmen. Bürgermeister wurde als\unsicher zusammenberuft
\end{verse}
\pstart
Wahlvorschlag in Seffern war      Stimmen, prop.\unsicher\ob
Kasbach und    Keil-Seffern mit    Stimmen.
\pend
